\begin{graphicspathcontext}{{./chapters/ai-fr/imgs/},{./chapters/ai-fr/imgs/auto/},\old}
	
	\begin{frame}[fragile,t]{{L'IA en santé :} Une histoire riche}
		L'IA est utilisée depuis des années dans le domaine de la santé, notamment en imagerie (détection de cancers) et à des fins organisationnelles (établissement de plannings)
	\end{frame}
	
	\begin{frame}[fragile,t]{{L'IA en santé :} De nouvelles perspectives}
		L'essor de l'IA générative et des LLMs ouvre aujourd'hui de nouvelles voies :
		\begin{itemize}
			\item Assistance au diagnostic
			\item Utilisation thérapeutique des agents conversationnels
			\item Formulation de contraintes de planning en langage naturel
		\end{itemize}
	\end{frame}
	
	\begin{frame}[fragile,t]{{L'IA en santé :} Assistance au diagnostic}
		% TODO
	\end{frame}
	
	\begin{frame}[fragile,t]{{L'IA en santé :} Utilisation thérapeutique des LLMs}
		% TODO Parler de l'étude de 2024 sur l'arrêt du tabac
	\end{frame}
	
	\begin{frame}[fragile,t]{{L'IA en santé :} Génération de plannings}
		% TODO
	\end{frame}
	
	\begin{frame}[fragile,t]{{L'IA en santé :} Des défis à résoudre}
		L'utilisation de l'IA en santé fait face à plusieurs défis :
		\begin{itemize}
			\item Faible niveau de preuve, peu d'études cliniques
			\item Le domaine de la santé est à haut risque
		\end{itemize}
	\end{frame}
	
\end{graphicspathcontext}

\endinput